\begin{minipage}{0.46\linewidth}
    Un skieur, de masse $M=\qty{80}{\kg}$ (avec son équipement), se trouve au
    sommet (point $O$) d'une piste rectiligne inclinée d'un angle
    $\alpha=\ang{4.59}$. La longueur de la pente est $L=\qty{200}{\m}$.

    \begin{donnees}
        \begin{itemize}
            \item Intensité de la pesanteur~: $g=\qty{9.8}{\N\per\kg}$.
        \end{itemize}
    \end{donnees}
\end{minipage}
\hfill
\begin{minipage}{0.53\linewidth}
\begin{figure}[H]
    \centering
    \begin{tikzpicture}[thick]
        \def\al{9}
        \def\b{-7}
        \node[minimum width=8cm,minimum height=2mm,top color=gray,
            anchor=north east] (sol) at (0.5,0) {};
        \PlanIncline{P}{O}{Q}{\b}{\al}{$\alpha$}
        \foreach \p/\a in {O/left,P/right}
            \node[point,label={[\a,inner sep=2pt]:$\p$}] at (\p) {};
        \node[anchor=south,font=\LARGE,rotate=18,yshift=-2.4mm] (S) at
            ($(O.center)!0.5!(P.center)$) {\faSkiing};
        \draw (Q) |- +(0.2,0.2) |- cycle;
        \draw[<->] ([shift=(90-\al:0.8)]O) coordinate (O') --
            node[above] {$L$} ([shift=(90-\al:0.8)]P) coordinate (P');
        \draw[dashed] (O) -- (O');
        \draw[dashed] (P) -- (P');
        \node[signal,draw,fill=black!10,anchor=west,solid,inner sep=2pt,
            signal to=west,xshift=5mm] at (P) {$E_{pp}=0$};
    \end{tikzpicture}
\end{figure}
\end{minipage}
\begin{enumerate}
    \item Le skieur part du sommet de la piste avec une vitesse de valeur
          $v_{O}=\qty{0.5}{\m\per\s}$.
    
          Calculer, au moment du départ, son énergie cinétique $E_{c}(O)$ et son
          énergie potentielle de pesanteur $E_{pp}(O)$.~\textbf{(2pt)}
    \item En appliquant la conservation de la somme $E_{c}+E_{pp}$, calculer
          l'énergie cinétique du skieur lorsqu'il est parvenu au
          bas de la piste (point $P$).~\textbf{(1pt)}

          En déduire la valeur de sa vitesse au point $P$.~\textbf{(1pt)}
\end{enumerate}

